% Generated draft: requires mathematical and source review.
\documentclass[11pt]{article}
\usepackage[margin=25mm]{geometry}
\usepackage{fontspec}
\setmainfont{Latin Modern Roman}
\usepackage{amsmath,amssymb,mathtools,unicode-math}
\setmathfont{Latin Modern Math}
\usepackage{enumitem,needspace,xurl,hyperref}
\setlength{\parindent}{0pt}
\setlength{\parskip}{5pt}
\setlength{\emergencystretch}{3em}
\newcommand{\N}{\mathbb{N}}
\newcommand{\Z}{\mathbb{Z}}
\newcommand{\Q}{\mathbb{Q}}
\newcommand{\R}{\mathbb{R}}
\newcommand{\C}{\mathbb{C}}
\newcommand{\F}{\mathbb{F}}
\newcommand{\A}{\mathbb{A}}
\newcommand{\PP}{\mathbb{P}}
\newcommand{\E}{\mathbb{E}}
\newcommand{\eps}{\varepsilon}
\newcommand{\dd}{\,\mathrm{d}}
\begin{document}
\section*{UnsolvedMath ID: 8700050}
\textbf{Source problem code:} catalogue entry
\subsection{Definitions and mathematical statement}
Let \(\mathbb G_m^n\) denote the algebraic torus over \(\mathbb Q\), with geometric points \((\overline{\mathbb Q}^{\times})^n\). A torsion point has every coordinate a root of unity, and a torsion translate is a coset \(\zeta H\) of an algebraic subgroup \(H\) by such a point.

For every integer \(n\ge1\), the source conjectures a constant \(c(n)>0\) such that every algebraic subvariety \(V\subset\mathbb G_m^n\), defined over \(\mathbb Q\) and irreducible over \(\mathbb Q\), which is not a union of torsion translates of algebraic subgroups, satisfies
\[
 \widehat h(V)\ge c(n)\deg(V)^{\frac{s-\dim V-1}{s-\dim V}},
\]
where \(s\) is the dimension of the smallest algebraic subgroup containing \(V\). The denominator is the positive dimension gap in the intended nonspecial situation. The constant depends only on \(n\), not on \(V\), its degree, dimension, or containing subgroup. No small-exponent loss or additional constant dependence is inserted.

Irreducibility over \(\mathbb Q\) is retained and is not replaced by geometric irreducibility. Likewise ``not a union of torsion translates'' is not replaced by the stronger, differently worded requirement of not lying in a specified proper torsion translate; the subgroup dimension \(s\) is included precisely as supplied.

The excerpt uses degree and canonical subvariety height without defining their embedding or normalization. A conventional degree uses the projective closure under \((x_1,\ldots,x_n)\mapsto[1:x_1:\cdots:x_n]\). The precise source normalization of \(\widehat h(V)\) must be recovered: total cycle height, height divided by degree, and essential minimum are different quantities and yield different degree exponents. The summary's ``essential/minimal'' description is therefore not silently substituted for the displayed height.

The supplied summary attributes qualitative/effective progress to Bombieri--Masser--Zannier and Amoroso--David, while saying the exact exponent remains open and the known statements use differing definitions or exponents. No precise theorem citations are supplied. The 2004 conjecture and later partial-progress assessment are source reports pending current primary-literature review.
\subsection{Short English statement}
Bound the canonical height of every eligible \(\mathbb Q\)-irreducible torus subvariety below by the stated power of its degree, uniformly in ambient dimension. The exact height normalization is essential, and the source's nonspecial and irreducibility hypotheses must be retained.
\subsection{Sources}
\begin{itemize}
\item Local registry: \url{data/unsolvedmath/problems.json}, permanent UnsolvedMath ID \(8700050\), source code catalogue entry; supplied statement, background, and research summary, updated 2026-08-06.
\item Michel Waldschmidt, \emph{Open Diophantine Problems} (2004), Conjecture 4.5, attributed to Amoroso--David: \url{https://arxiv.org/abs/math/0312440}. Supplied source-TeX extraction accessed 2026-07-29; a current primary-literature status audit remains required.
\item Bombieri--Masser--Zannier and Amoroso--David are named by the summary for related lower bounds, without separate precise publication details or theorem statements.
\end{itemize}
\end{document}
